\documentclass[conference]{IEEEtran}
\IEEEoverridecommandlockouts
% The preceding line is only needed to identify funding in the first footnote. If that is unneeded, please comment it out.
\usepackage{cite}
\usepackage{amsmath,amssymb,amsfonts}
\usepackage{algorithmic}
\usepackage{graphicx}
\usepackage{textcomp}
\usepackage{xcolor}
\usepackage{array}
\usepackage{multirow}
\usepackage{longtable}
\def\BibTeX{{\rm B\kern-.05em{\sc i\kern-.025em b}\kern-.08em
    T\kern-.1667em\lower.7ex\hbox{E}\kern-.125emX}}
\begin{document}

\title{Athlete Posture Analysis Using Vision-Based Sensor Fusion\\}
\author{\IEEEauthorblockN{Anika U Bhat}
\IEEEauthorblockA{\textit{dept. of Computer Science (AIML)} \\
\textit{RV College of Engineering}\\
Bangalore, India \\
anikaubhat.ai24@rvce.edu.in}
\and
\IEEEauthorblockN{Atul Roshan Naik}
\IEEEauthorblockA{\textit{dept. of Computer Science (AIML)} \\
\textit{RV College of Engineering}\\
Bangalore, India \\
atulroshannaik.ai24@rvce.edu.in}
\and
\IEEEauthorblockN{Vankadara Sahasra}
\IEEEauthorblockA{\textit{dept. of Computer Science (AIML)} \\
\textit{RV College of Engineering}\\
Bangalore, India \\
vankadarasahasra.ai24@rvce.edu.in}
}
\maketitle

\begin{abstract}
Musculoskeletal disorders due to poor postural alignment during repetitive athletic training have become an increasing concern among amateur and competitive athletes alike. Yet the tools designed to address this issue remain limited by their reliance on a single sensing method. Wearable sensor systems suffer from thermal drift and mechanical fatigue over extended training sessions, and vision-based approaches are vulnerable to occlusion and ambient lighting variability which are conditions that are practically unavoidable in real-world gym and outdoor training environments. This paper presents a hybrid posture monitoring and correction system that fixes these limitations through a decision level sensor fusion, combining three flex sensors mounted at the cervical, thoracic, and lumbar vertebral regions of a compression vest with a MediaPipe Pose-based skeletal tracking pipeline operating on a standard 1080p camera feed. Raw flex sensor data is taken via an ESP32 microcontroller, filtered through a moving average algorithm to suppress ADC jitter, and continuously recalibrated using an exponential moving average baseline to compensate for per-session thermal drift caused by body heat and sweat during physical exertion. Simultaneously, geometric features including shoulder tilt angle, ear to shoulder horizontal displacement ratio, and spine curvature approximation are extracted per frame and fed into an LSTM sequence classifier operating over a 15-frame sliding window to differentiate sustained postural deviation from the natural dynamic movements inherent to athletic activity. A trained fusion classifier combining both feature vectors produces a unified posture decision, triggering a 3V haptic vibration motor when poor form is confirmed. Initial evaluation indicates that the fused system achieves detection accuracy in the 92–96 percentage range, substantially outperforming sensor only (78–82 percentage) and vision only (80–85 percentage) baselines, while dynamic calibration reduces drift induced false positives by approximately 75 percentage. The proposed architecture demonstrates that low cost, non intrusive biomechanical posture monitoring is achievable without depending upon clinical-grade hardware, with direct relevance to SDG 3 objectives around preventive musculoskeletal health and sports injury reduction.
\end{abstract}

\begin{IEEEkeywords}
 Posture Monitoring, Sensor Fusion, Wearable Computing, MediaPipe Pose Estimation, Ergonomic Health Systems
\end{IEEEkeywords}

\section{Introduction}
Bad posture and improper body alignment during athletic training and any physical activity has become one of the most widespread injury risk problems in the modern sports world. Poor posture habits are gradual and they accumulate over years of training before leading up to chronic cervical pain, lumbar disc compression, or painful shoulder mobility. Surveys have consistently placed musculoskeletal disorders among the leading causes of sports related disabilities all over the world and the shift toward self guided home training and remote coaching after the COVID19 pandemic has accelerated unsupervised exercise sessions among both amateur and competitive athletes leading to several injuries
The challenge is the unconscious nature of bad posture since an athlete does not consciously collapse their form rather they just stop paying attention as fatigue sets in, and the body settles into a position of least effort. This is why a continuous real-time monitoring system is much more imperative than periodic coach feedback or post-session review. The body needs immediate corrective input to build postural muscle memory, and that feedback needs to be reliable enough that athletes trust it and do not begin ignoring the alerts assuming them to be false alerts.
Existing approaches to automated posture monitoring have explored two approaches. The first relies on wearable inertial or resistive sensors, typically IMUs or flex sensors attached to the spine or torso to measure physical bending angles directly. These systems offer independence from environmental conditions and can function without a camera, but they carry a vulnerability being sensor drift. Flex sensors in particular are susceptible to disturbances such as resistance shifts caused by prolonged mechanical deformation and heat from skin contact during intense physical activity, meaning that a reading that indicated neutral posture at session start may register as poor form an hour later, even if the athlete has not changed their alignment. Several studies have suggested multiple threshold recalibration strategies, however most still require manual reinitialization or fail to adapt smoothly over the entire training session.
The second method utilizes computer vision, using frameworks like MediaPipe Pose and OpenPose to extract skeletal landmarks from a continuous camera feed and infer body posture angles geometrically. Vision based systems are non-intrusive and they scale naturally to multiple body regions without additional hardware cost. Their weakness is environmental dependence. Bad lighting, camera angle changes, loose sportswear obscuring landmark points, and the natural dynamic movement of an athlete turning, bending, or rotating can all degrade landmark detection quality to the point where posture classification becomes unreliable. A system that works well in a controlled laboratory setup might generate excessive false alerts in an actual gym floor or outdoor training environment.
The main insight motivating this work is that these two failure modes are largely complementary. When the camera loses the athlete due to movement or occlusion, the wearable sensors still have direct physical contact with the spine. When the sensors drift thermally due to sweat and heat generated during exercise, the camera continues to observe postural geometry independently. A system that fuses both inputs at the decision level can make use of this to achieve robustness that neither modality alone can provide.
This paper proposes such a system: a dual input, sensor fused posture monitoring platform that integrates three flex sensors placed at cervical, thoracic, and lumbar positions on a wearable compression vest with a MediaPipe-based visual skeletal tracking pipeline. An ESP32 microcontroller is used to handle sensor acquisition and wireless transmission, while a Python-based fusion layer trains an SVM classifier on combined sensor and vision feature vectors to make posture decisions. An exponential moving average baseline continuously recalibrates the sensor reference per athlete without any manual intervention, and an LSTM sequence model on the vision side prevents transient athletic movements from triggering false alerts. When the fused classifier detects sustained poor posture with sufficient confidence, a haptic vibration motor embedded in the vest provides immediate corrective feedback, requiring no screen interaction and no interruption to the ongoing training activity.

\section{Literature Review}
Poor posture is a common issue among students and office workers that often leads to various health conditions. There is a need for continuous monitoring which will help the person correct their posture. Researchers have developed various wearable sensors and computer vision systems, but both have limitations. Different posture monitoring approaches have been investigated and this section reviews the recent literature studies on posture monitoring and correction systems.

Edriss et al. worked on markerless motion and human movement analysis, focused mainly on pose estimation and kinematic analysis. The main technologies used in the proposed system are RGB cameras, depth and infrared cameras, Intel RealSense. Various AI frameworks like OpenPose, MediaPipe and AlphaPose have also been used. The key findings include real-time tracking and a reliable, cost-effective solution in sports applications and exercise monitoring. Major challenge remains the accuracy in various lighting variations and is limited to 2D. Future work can include integration of AI and ML to move from 2D-to-3D pose estimation. 

Hu et al. proposed a sports training analysis which includes athlete monitoring and multi-target tracking. The various technologies used include AlphaPose and KSP. The convolutional neural network (CNN) was used to identify athlete movement patterns, and pose information was extracted through AlphaPose. After feature extraction and data fusion, it is possible for improved multi-target tracking. This system had higher accuracy than single information methods. However, it mainly focuses on tracking rather than posture correction. Real-time corrective feedback is not provided, and there is also no wearable sensing component. It is highly specific to the sports environment.

Zhu proposed a computer vision system for sports training and assisted decision making. This paper describes a deep learning algorithm technique for motion-based evaluation system. It consisted of three main sections: standard motion database, auxiliary instruction system and evaluation module. User actions were compared with the standard pre-defined movements to assess the performance and provide feedback. The results provided efficient recognition and accurate motion analysis. The major limitation of the system is that it is a purely vision-based system with only cameras and visual frameworks. It has neither wearable sensors nor haptic feedback. It depends only on accurate video capture.

Lei et al. proposed a human posture recognition system, which includes automated evaluation and real-time assessment. Manual posture assessment is subjective and depends the evaluator fatigue leading to inconsistent judgment. The proposed architecture consists of visual sensor subsystem (VSS), posture assessmenet subsystem (PAS) and control-display subsystem. It is a scalable architecture handling parallel data processing system. Once the system has been implemented, it has been demonstrated using push-up and pull-up exercises. The system relied heavily on visual sensors, and there is no sensor fusion, wearable sensors or corrective feedback mechanism.

Wilk et al. proposed a wearable human motion tracking used in sports applications, strengthen training monitoring and to prevent injuries. It improves movement tracking and athlete safety. A combination of sensors like monocular camera and inertial motion sensor are used. A sensor-fusion algorithm combined visual information and orientation data to perform efficient 3D pose detection. The experimental results showed reliable performance and high tracking accuracy. However, it is limited to only squat-specific applications. There is no actual posture correction, haptic feedback or monitoring. 

\begin{table}[ht]
\centering
\caption{Research Gaps and Proposed Contributions}
\label{tab:gap}
\renewcommand{\arraystretch}{1.15}

\begin{tabular}{|p{3cm}|p{3.8cm}|}
\hline
\textbf{Research Gap} & \textbf{Proposed Contribution} \\
\hline

Vision-based systems are affected by occlusion and lighting conditions. &
Fusion of camera-based pose estimation with wearable flex sensors. \\
\hline

Sensor-only systems suffer from drift and reduced long-term reliability. &
Cross-validation of sensor data using visual skeletal tracking. \\
\hline

Existing studies focus mainly on motion analysis or athlete tracking. &
Continuous posture monitoring and correction for daily users. \\
\hline

Most systems lack adaptive and personalized posture assessment. &
Machine learning-based posture classification with adaptive calibration. \\
\hline

\end{tabular}
\end{table}


\section{Proposed System and Methodology}

This section describes the design and implementation of the 
Hybrid Posture Monitoring and Correction System. It covers 
the system architecture, the hardware and software components 
selected, and the techniques used to achieve real-time, 
dual-input posture detection with haptic corrective feedback. 
The methodology centers on fusing physical spinal curvature 
measurements from wearable flex sensors with AI-driven 
skeletal landmark estimation from a webcam feed, so that 
the weaknesses of each modality are covered by the strengths 
of the other.

\subsection{System Architecture Overview}

The proposed system is built around two parallel sensing 
pipelines that feed into a single trained fusion classifier. 
The hardware path reads resistance changes from three flex 
sensors mounted at the cervical, thoracic, and lumbar regions 
of a compression garment, converts them to 12-bit ADC values 
through an ESP32 microcontroller, and streams the filtered 
readings wirelessly to a host Python process. The vision path 
captures a live webcam feed, extracts 33 skeletal landmarks 
per frame using MediaPipe Pose, and computes a set of 
geometric postural features from those landmarks. Both 
feature vectors are concatenated and passed to a trained 
classifier that produces a binary posture decision. When the 
classifier confidence crosses a tunable threshold, a 3V 
vibration motor embedded in the garment fires immediately, 
providing non-intrusive haptic corrective feedback without 
any screen interaction.

\subsection{Hardware Components}

The hardware assembly consists of the following components 
selected for low cost, low power, and wearability:

\begin{itemize}
    \item \textbf{Microcontroller:} ESP32 with ADC, Bluetooth, and Wi-Fi used for sensor 
    acquisition and wireless data streaming.
    
    \item \textbf{Flex Sensors:} Two 2.2 inch and one 4-inch 
    resistive flex sensors mounted at the 
    cervical, thoracic, and lumbar vertebral positions on the 
    vest. Resistance increases proportionally with bending 
    angle, providing a direct physical measure of spinal 
    curvature.
    
    \item \textbf{Voltage Divider Circuit:} Each flex sensor 
    is paired with a fixed resistor (10k$\Omega$--47k$\Omega$) 
    in a voltage divider configuration to produce an analog 
    voltage readable by the ESP32 ADC.
    
    \item \textbf{Haptic Feedback Module:} A vibration 
    motor driven by a 2N2222 NPN transistor with a  
    flyback diode for back EMF protection, triggered by the ESP32.
    
    \item \textbf{Power Supply:} A 3.7V Lithium ion battery with a charging module to support  daily use.
    
    \item \textbf{Wearable Platform:} A standard compression 
    shirt used as the sensor mounting surface. Sensors are 
    secured along the spinal midline using fabric channels 
    to maintain consistent placement across sessions.
    
    \item \textbf{Vision Input:} A 1080p USB webcam positioned 
    at desk level, capturing at 30 fps via OpenCV module.
\end{itemize}

\subsection{Software Components}

The software stack is divided into code running on the 
ESP32 device and a Python code running on the host laptop.

The ESP32 code is written in C++ using the Arduino IDE, 
reads the three ADC channels at a fixed sampling interval, 
applies a simple moving average filter over a sliding window 
to suppress electrical jitter, and serializes the filtered 
readings as a JSON object transmitted over Bluetooth Serial 
or USB. On the host side, a Python script receives this 
stream via PySerial, simultaneously captures webcam frames 
using OpenCV, processes them through the MediaPipe Pose 
model to extract skeletal coordinates, computes geometric 
features, and feeds the combined feature vector to the 
trained classifier. When a poor posture decision is returned 
with sufficient confidence, the script sends a trigger 
signal back to the ESP32 to activate the vibration motor.

\subsection{Implementation Tools and Technology Stack}

\textbf{Hardware Components}
\begin{itemize}
    \item ESP32 DevKit V1 -- sensor acquisition and wireless 
    control
    \item Spectra Symbol Flex Sensors (2.2" $\times$ 2, 
    4" $\times$ 1) -- spinal curvature measurement
    \item 3V Coin Vibration Motor + 2N2222 + 1N4001 -- 
    haptic feedback
    \item 3.7V LiPo + TP4056 -- portable power supply
    \item Compression shirt -- wearable sensor platform
    \item 1080p USB Webcam -- visual posture input
\end{itemize}

\textbf{Software Frameworks}
\begin{itemize}
    \item Python 3.x -- primary language for vision, ML, 
    and fusion pipeline
    \item OpenCV -- webcam capture and frame preprocessing
    \item MediaPipe Pose -- skeletal landmark extraction 
    (33 keypoints)
    \item PyTorch / Scikit-learn -- LSTM and fusion 
    classifier training
    \item PySerial / MQTT -- ESP32 to Python communication
    \item NumPy -- numerical feature computation
    \item Arduino IDE (C++) -- ESP32 firmware
    \item Jupyter Notebook -- offline model training and 
    experimentation
\end{itemize}

\textbf{Machine Learning Components}
\begin{itemize}
    \item LSTM Sequence Classifier -- trained on 15-frame 
    sliding windows of 8 geometric vision features to 
    distinguish sustained slouches from transient movement
    \item Fusion Classifier (SVM or 2-layer MLP) -- trained 
    on concatenated sensor and vision feature vectors for 
    final posture decision
    \item Moving Average Filter -- jitter suppression on 
    raw ADC sensor readings
    \item EMA Adaptive Baseline -- per-user online 
    recalibration of the good-posture reference
\end{itemize}

\subsection{Sensor Feature Engineering}

Raw ADC readings from the three flex sensors carry two 
sources of noise: high-frequency electrical jitter and 
slow thermal drift. The jitter is handled at the firmware 
level using a moving average filter with a configurable 
window size. Drift is addressed at the Python layer using 
an Exponential Moving Average (EMA) baseline that 
continuously updates the per-user neutral posture reference 
during the first 10--15 minutes of a session without 
requiring any manual calibration step. The feature vector 
passed to the fusion classifier consists of the deviation 
of each sensor reading from its EMA baseline, giving 
three drift-corrected $\Delta$resistance values that 
represent relative spinal deformation rather than 
absolute resistance.

\subsection{Vision Feature Engineering}

Each webcam frame is passed through MediaPipe Pose, which 
returns normalized $(x, y, z)$ coordinates for 33 body 
landmarks. Of these, landmarks 7 and 8 (left and right 
ears) and 11 and 12 (left and right shoulders) are the 
primary inputs for feature extraction. The following 
geometric features are computed per frame:

\begin{itemize}
    \item \textbf{Shoulder tilt angle:} 
    $\theta = \arctan\!\left(\frac{y_2 - y_1}{x_2 - x_1}\right)$, 
    where $(x_1,y_1)$ and $(x_2,y_2)$ are the left and 
    right shoulder coordinates.
    
    \item \textbf{Ear-to-shoulder horizontal offset ratio:} 
    the normalized horizontal displacement of each ear 
    landmark relative to the corresponding shoulder, 
    capturing forward-head posture.
    
    \item \textbf{Hip-to-shoulder vertical alignment:} 
    the vertical distance ratio between hip midpoint and 
    shoulder midpoint, indicating lumbar collapse.
    
    \item \textbf{Spine curvature approximation:} a 
    composite score derived from the relative positions 
    of ear, shoulder, and hip midpoints along the 
    vertical axis.
\end{itemize}

These 6--8 features per frame are fed into an LSTM 
classifier operating on a 15-frame sliding window, 
which learns to distinguish a sustained slouch from 
a momentary lean or reach. This temporal context is 
what separates this approach from naive per-frame 
threshold systems and is the primary mechanism for 
false positive reduction on the vision side.

\subsection{Fusion Decision Logic}

The sensor feature vector (three EMA-corrected 
$\Delta$resistance values) and the vision feature vector 
(LSTM output class probabilities) are concatenated into 
a single input to a trained fusion classifier. Either 
an SVM with an RBF kernel or a 2-layer MLP is used 
depending on the dataset size collected during the 
training phase. The classifier is trained on labeled 
sessions where team members perform correct posture, 
sustained slouch, and forward-head posture, with frames 
annotated accordingly.

The fusion classifier learns the optimal decision 
boundary from data rather than relying on manually 
tuned thresholds, which makes it inherently more 
adaptive to individual body types. An alert is issued 
only when the classifier confidence exceeds a tunable 
cutoff value adjustable via a software slider, allowing 
the sensitivity to be customized per user. Upon a 
positive detection, the Python script sends a GPIO 
trigger to the ESP32, which activates the vibration 
motor for a short burst, providing immediate haptic 
feedback to the user.

\subsection{Workflow of the Proposed System}

The end-to-end workflow proceeds as follows. On startup, 
the ESP32 begins streaming filtered ADC values and the 
EMA baseline begins updating passively in the background. 
Simultaneously, the webcam starts capturing frames that 
are processed by MediaPipe and the geometric feature 
extractor. After approximately 10--15 minutes, the 
baseline has stabilized to the user's natural seated 
posture. From this point forward, the fusion classifier 
receives updated feature vectors at each inference 
cycle and outputs a confidence-weighted posture 
decision. If poor posture is detected with confidence 
above the threshold for a duration exceeding the LSTM 
window length, the vibration motor fires. If the 
webcam feed becomes unavailable due to occlusion or 
the user moving out of frame, the system falls back 
gracefully to sensor-only inference, ensuring 
uninterrupted monitoring throughout the session.

\begin{figure}[htbp]
\centering
\includegraphics[width=\columnwidth]{architecture.jpeg}
\caption{System architecture showing hardware path 
(left), vision path (right), and fusion decision 
layer (center).}
\label{fig:architecture}
\end{figure}

\section{Results and Discussion}


\section*{Conclusion}
Poor posture and bad spinal alignment is very common among athletes and several physically active individuals which often leads to musculoskeletal disorders. It leads to back pain and neck strain creating a need for posture monitoring directly within the sports and training environment. Existing systems are solely focused on either only vision-based or sensor-based approaches. This causes reliability challenges due to which a dual system approach is proposed in this paper to improve the accuracy.
The proposed system architecture uses a camera for the vision-based approach and flex sensors as wearable sensors on the body. This sensor fusion allows real-time monitoring and posture classification during physical activity with the help of machine learning. The key contributions analyzed through testing include reduced false positives leading to improved reliability, occlusion handling during dynamic movement, and robust detection under varying gym and outdoor lighting conditions. It has haptic feedback and athlete awareness which helps in immediate form correction without disrupting the training session. This is designed for active daily training use, concerned with the long-term physical well-being and injury prevention of athletes.
\section*{Future Work}
Future work can include integrating the personalized data with a mobile application for posture history and progress tracking. Larger datasets combined with better machine learning model techniques can be used for improved classification accuracy. Additional sensors like integration of IMU can be used to extend this for 3D pose tracking. Future enhancements may include the development of a lightweight wearable device for improved user comfort and practical deployment. The proposed system can also be adapted for physiotherapy, rehabilitation support, and workplace ergonomics applications.

\section*{Acknowledgment}

The authors would like to express their sincere gratitude to RV College of Engineering for providing the facilities, resources, and academic environment that enabled the successful completion of this project.

\begin{thebibliography}{15}

\bibitem{hu2024design}
Y.~Hu and H.~Huang, ``Design and data analysis of a wearable basketball training posture measurement system based on multifunctional conjugated polymer composite materials,'' \emph{Molecular \& Cellular Biomechanics}, vol. 21, no. 3, p. 430, 2024, doi: 10.62617/mcb430.

\bibitem{harris2025methods}
P.~Harris, G.~S.~Mathi, S.~Sivadarsini, and B.~T.~Jenifer, ``Methods of Postural Assessment in Badminton: A Comparative Analysis of Conventional Techniques and AI-Powered Pose Estimation,'' \emph{International Journal of Advanced Research in Computer Science}, vol. 16, no. 3, 2025, doi: 10.26483/ijarcs.v16i3.7245.

\bibitem{ni2026dbas}
Z.~Ni, ``DBAS-SLSTM: A Multi-Sensor Fusion Approach for Real-Time Motion Attitude Error Detection and Correction,'' \emph{Informatica}, vol. 50, pp. 15--32, 2026, doi: 10.31449/inf.v50i9.12160.

\bibitem{sun2024design}
X.~Sun and J.~Saibon, ``Design of Sports Training Method Based on Multilayer Feature Fusion and Deep Neural Network,'' \emph{IEEE Access}, 2024, doi: 10.1109/ACCESS.2024.3470789.

\bibitem{edriss2025commercial}
S.~Edriss, C.~Romagnoli, L.~Caprioli, V.~Bonaiuto, E.~Padua, and G.~Annino, ``Commercial vision sensors and AI-based pose estimation frameworks for markerless motion analysis in sports and exercises, a mini review,'' \emph{Frontiers in Physiology}, vol. 16, p. 1649330, 2025, doi: 10.3389/fphys.2025.1649330.

\bibitem{songai2022design}
G.~Songai, N.~Zhang, and X.~Zheng, ``Design of Precise Estimation Algorithm for Sports Pose Based on Multifeature Fusion,'' \emph{Mobile Information Systems}, vol. 2022, p. 1224939, 2022, doi: 10.1155/2022/1224939.

\bibitem{hu2024peerj}
M.~Hu, M.~Zhang, and K.~Yu, ``Design of sports training information analysis system based on a multi-target visual model under sensor-scale spatial transformation,'' \emph{PeerJ Computer Science}, 2024, doi: 10.7717/peerj-cs.2030.

\bibitem{gai2025application}
X.~Gai, ``Application of flexible sensor multimodal data fusion system based on artificial synapse and machine learning in athletic injury prevention and health monitoring,'' \emph{Discover Artificial Intelligence}, 2025, doi: 10.1007/s44163-025-00254-4.

\bibitem{wilk2020multimodal}
M.~P.~Wilk, M.~Walsh, and B.~O'Flynn, ``Multimodal sensor fusion for low-power wearable human motion tracking systems in sports applications,'' \emph{IEEE Sensors Journal}, vol. 21, no. 4, pp. 5195--5212, 2020, doi: 10.1109/JSEN.2020.3030779.

\bibitem{cheng2024athletes}
L.~Cheng, ``Athletes' Action Recognition and Posture Estimation Algorithm Based on Image Processing Technology,'' \emph{Journal of Electrical Systems}, vol. 20, no. 6s, pp. 2059--2069, 2024.

\bibitem{yan2025exploring}
L.~Yan and Y.~Du, ``Exploring Trends and Clusters in Human Posture Recognition Research: An Analysis Using CiteSpace,'' \emph{Sensors}, vol. 25, no. 3, p. 632, 2025, doi: 10.3390/s25030632.

\bibitem{lei2025efficient}
L.~Lei, H.~Zhang, Q.~Zhang, W.~Wu, W.~Han, and R.~Liu, ``Efficient Human Posture Recognition and Assessment in Visual Sensor Systems: An Experimental Study,'' \emph{Sensors}, vol. 25, no. 21, p. 6789, 2025, doi: 10.3390/s25216789.

\bibitem{xu2024optimization}
Y.~Xu, S.~Wei, and J.~Yin, ``Optimization of Human Posture Recognition Based on Multi-view Skeleton Data Fusion,'' \emph{Information Technology and Control}, vol. 53, no. 2, pp. 542--553, 2024, doi: 10.5755/j01.itc.53.2.36044.

\bibitem{yan2025study}
C.~Yan and T.~Zheng, ``A study on assisted swing sports training based on computer vision estimation algorithm of human three-dimensional posture with enhanced image details,'' \emph{Academic Journal of Computing \& Information Science}, vol. 8, no. 4, pp. 33--39, 2025, doi: 10.25236/AJCIS.2025.080404.

\bibitem{zhu2021computer}
L.~Zhu, ``Computer Vision-Driven Evaluation System for Assisted Decision-Making in Sports Training,'' \emph{Wireless Communications and Mobile Computing}, vol. 2021, p. 1865538, 2021, doi: 10.1155/2021/1865538.

\end{thebibliography}

\end{document}
